\sourcepage{312}
\chapterappendix{10A}{用有向无环图表示变量间关联}
\sourcepage{313}
使用\term{有向无环图}（directed acyclic graphs，DAGs，英文发音与rags押韵），改进了根据观察性资料建立变量关联模型的方法。DAG以示意图表示研究者认为各研究变量如何相互联系。DAG背后的数学原则，使因果效应的估计成为可能。

DAG用单向箭头表示两个变量间的因果方向。因此，$A\to B$表示研究者认为预测因素$A$引起结局$B$，例如过量饮酒引起髋部骨折。它被称为有向无环图，是因为变量构成的网络以有方向的箭头连接、不含环路，并用图形表示。

DAG可能连接十个或更多变量，因而变得复杂。但当只有三个变量，即研究者认为会引起某个结局的预测因素、该结局，以及可能与这两个变量有关的第三变量时，就相对容易理解。例如，研究低维生素D水平是否引起髋部骨折时，研究者可能担心户外运动的混杂作用，因为户外运动既能提高维生素D水平，也能改善力量和平衡。当然，研究者还可能关注跌倒史、营养状况和居住地等许多因素。不过，基本方法仍很直接：建立因果模型，也就是画出DAG以后，研究者必须决定哪些变量（如果有）会影响参与者的选择方式，以及除了预测因素和结局以外，还需测量哪些变量，才能在资料分析时将它们纳入统计模型。

与第1章介绍临床研究时一样，下面先讨论DAG的“解剖”，用图说明并命名它的各种形式，再讨论其“生理”，即DAG如何帮助进行因果推断和选择样本。

\subsection{DAG的例子}
在研究某个预测因素是否引起结局时，第三变量与二者只有八种联系形式（图10A.1—10A.8）。下面将逐一介绍这八种可能。预测变量始终用紫色文字表示，因为英文predictor与purple都以p开头，\sourcepage{314}结局变量则用赭色表示。请特别注意第三变量（青绿色）的位置，以及与其连接的蓝色箭头；各例之间的差别有时很细微。我们遵循惯例，让连接箭头指向所假定的因果方向，并尽可能从上到下、从左到右。

\subsection{未连接全部三个变量的DAG}
前两种情况，是第三变量只引起预测因素、不引起结局（图10A.1），或只引起结局、不引起预测因素（图10A.2）。在这两种DAG中，第三变量都与“预测因素是否引起结局”这一问题无关。
\par
\begin{figure}[htbp]\centering\chTenDAG{1}
\caption*{\textbf{图10A.1\quad 未连接全部三个变量的关联。}估计头部创伤对帕金森病的因果效应，不必考虑戴自行车头盔史；这也是以戴头盔为工具的工具变量分析DAG。}
\end{figure}
\par
\par
\begin{figure}[htbp]\centering\chTenDAG{2}
\caption*{\textbf{图10A.2\quad 未连接全部三个变量的关联。}男性不可能解释A型血与COVID-19住院的关联。}
\end{figure}
\par

例如，假设研究者想知道预测因素“严重头部创伤史”是否引起帕金森病（图10A.1）。第三变量“骑车时不戴头盔”会引起这一预测因素，因此有从戴自行车头盔指向头部创伤的箭头。但是，研究者并不认为戴或不戴头盔会直接引起帕金森病，因此二者之间没有箭头。这样就不必担心，戴自行车头盔会影响头部创伤对帕金森病因果效应的估计。

\sourcepage{315}再考虑A型血是否增加因COVID-19住院风险的研究（图10A.2）。男性因COVID-19住院的风险更高，因此有从男性性别指向COVID-19住院的箭头。\footnote{译注：底本此处标作图10A.3，相应图实际为图10A.2。}但是，性别与A型血没有关联，男女中A型血同样常见，因此没有从性别指向A型血的箭头。于是，不必担心男性性别会解释A型血对COVID-19住院的因果效应。尽管如此，在分析模型中加入性别，可能提高血型对住院影响的估计精度，也可能有助于考察性别是否修饰这一效应。

接下来的两种情况，是第三变量由预测因素引起，却与结局无关（图10A.3），或由结局引起，却与预测因素无关（图10A.4）。
\par
\begin{figure}[htbp]\centering\chTenDAG{3}
\caption*{\textbf{图10A.3\quad 未连接全部三个变量的关联。}男性型脱发不影响雄激素水平对肾癌的因果效应。}
\end{figure}
\par
\par
\begin{figure}[htbp]\centering\chTenDAG{4}
\caption*{\textbf{图10A.4\quad 未连接全部三个变量的关联。}使用镇痛药不改变铅暴露对关节病的因果效应估计。}
\end{figure}
\par

假设研究者想确定男性雄激素水平升高是否引起肾癌，并担心男性型脱发可能产生影响。虽然雄激素水平较高的男性可能更容易出现男性型脱发，却没有理由认为脱发引起肾癌。因此，不必调整脱发的影响，甚至无须询问这一情况（图10A.3）。

\sourcepage{316}另一种情况是，第三变量（如使用镇痛药）由结局（退行性关节病）引起，而不是由预测因素（铅暴露）引起；铅暴露只有在引起关节病的限度内，才会引起镇痛药使用（图10A.4）。同样，研究铅暴露与关节病的关联时，无须调整镇痛药使用。

\subsection{连接全部三个变量的DAG}
第三变量与预测因素和结局都相连，共有四种形式（图10A.5—10A.8）。这些情况需要额外注意，因为研究者必须决定分析资料时是否纳入第三变量，有时甚至在制定研究选择标准时就须作出决定。后面的“用DAG作因果推断”及“DAG与抽样”将讨论这一点。
\subsubsection{中介}
四种形式中，中介最容易理解（图10A.5）。它发生在预测因素（如血清LDL胆固醇水平）引起另一个因素（颈动脉狭窄），后者再引起结局（卒中）时。之所以称为中介，是因为LDL胆固醇对卒中的部分或全部影响，可能通过颈动脉狭窄传递。
\par
\begin{figure}[htbp]\centering\chTenDAG{5}
\caption*{\textbf{图10A.5\quad 中介。}LDL胆固醇对卒中的部分因果效应可能经颈动脉狭窄中介。}
\end{figure}
\par
\subsubsection{共同效应}
下一个DAG（图10A.6）表示\term{共同效应}（shared effect）：箭头相遇的点，也就是节点，常称为碰撞点。预测因素和结局都引起第三变量时，就会出现这种情况。以低热量饮食是否引起胃癌的研究为例，低热量饮食和胃癌都会使体重下降，因此体重下降就是共同效应。
\sourcepage{317}
\par
\begin{figure}[htbp]\centering\chTenDAG{6}
\caption*{\textbf{图10A.6\quad 共同效应。}按体重下降分层，可使低热量饮食对胃癌因果效应估计偏倚。}
\end{figure}
\par
共同效应可能造成问题。将碰撞点作为纳入条件，或将其纳入统计模型，可能引起偏倚，甚至使真实的因果关联消失。本例中，如果研究只纳入体重下降者，就会创造出低热量饮食与胃癌的负关联：只要二者之一存在，就无须另一个来解释体重下降。这可能使低热量饮食看似具有预防胃癌的作用，或掩盖二者之间真实的正关联。

\subsubsection{混杂}
最为人熟悉的DAG（图10A.7）表示混杂，即第三因素同时引起预测因素和结局。例如，研究吸大麻是否引起肺癌时，吸烟是潜在混杂因素，因为它既引起肺癌，也增加吸大麻的可能性。
\par
\begin{figure}[htbp]\centering\chTenDAG{7}
\caption*{\textbf{图10A.7\quad 混杂。}吸烟可促成吸大麻与肺癌的关联。}
\end{figure}
\par
如果吸大麻对肺癌的作用完全可由吸烟解释，在分析模型中调整吸烟后，吸大麻对肺癌的作用就会消失。也可能只有部分混杂，此时调整会使效应减弱。甚至还可能存在反向混杂，即抑制：调整混杂因素后，关联反而增强，见正文“低估因果效应”。

\subsubsection{并非DAG的“DAG”}
\sourcepage{318}
最后一种情况是因果图中存在环路（图10A.8）：预测因素（肥胖）引起结局（骨关节炎），结局又引起第三因素（久坐生活方式），第三因素再影响预测因素（肥胖）。
\par
\begin{figure}[htbp]\centering\chTenDAG{8}
\caption*{\textbf{图10A.8\quad 循环图。}DAG的原则不适用于此图。}
\end{figure}
\par
这种情况实际上不是DAG，因为它并非无环，因此DAG的规则不适用。处理这种情况需要特殊方法，例如随时间反复测量这些变量，以解开哪个是因、哪个是果这一“先有鸡还是先有蛋”的问题。

\subsection{用DAG作因果推断}
画出DAG并判断它属于上述哪些情况后，下一步必须决定是否处理第三变量或变量集，以及如何处理。答案取决于图中呈现的是哪种情况。

DAG指导研究者如何设计研究、分析资料，以估计预测因素对结局的因果效应。可以把箭头想象为管道或路径，因果信息沿着它从一个变量流向另一个变量。不论箭头指向哪边，管道都允许信息流动，因此一般称为“开放”路径。

有一个重要例外：两个不同变量的箭头指向同一个变量时，这两个变量间的路径称为“关闭”。例如，图10A.6中，从低热量饮食经体重下降到胃癌的路径就是关闭的。记忆这一规则时，可以想象彼此相向的箭头挡住了对方。使用DAG的经验越多，就越容易识别哪些路径开放、哪些关闭，并理解为什么这很重要。可公开使用的软件，如\url{daggity.net}，能提供帮助。

规划资料分析时，模型中应纳入哪些变量？将某个变量纳入分析模型，就等于调整该变量，这一过程也称为对该变量“\term{条件化}”（conditioning）。\sourcepage{319}首先，应该包含从预测因素到结局的直接路径（预测因素$\to$结局），因为这正是要估计的因果效应。因此，模型中总要包括预测因素和结局。

但是，我们不希望涉及共同效应或混杂的路径有“因果流”（表10A.1）。幸运的是，可以利用DAG的一条重要规则来防止它：\emph{调整开放路径上的变量，会关闭该路径；调整使路径关闭的变量，则会打开该路径。}具体来说，应在调整模型中纳入每条混杂路径上的某个变量，将所有混杂路径关闭。相比之下，共同效应路径应当保持原样，因为它已经关闭，我们不希望把它打开。

\begin{longtable}{L{11mm}L{34mm}L{14mm}L{14mm}L{32mm}L{36mm}}
\caption*{\textbf{表10A.1\quad 使路径保持开放或关闭的指导}}\\\toprule
情形&例图&蓝箭头路径当前状态&蓝箭头路径期望状态&行动&这一行动回答的问题\\\midrule\endfirsthead
\toprule 情形&例图&当前状态&期望状态&行动&回答的问题\\\midrule\endhead
中介&\chTenTablePath{1}&开放&保持开放&不在模型中纳入颈动脉狭窄，以确定LDL胆固醇的总效应&LDL胆固醇对卒中风险的因果效应，包括它对颈动脉狭窄的影响，是多少？\\
中介&\chTenTablePath{1}&开放&关闭路径&在模型中调整颈动脉狭窄，以确定LDL胆固醇的直接效应$^{a}$&LDL胆固醇不通过颈动脉狭窄发挥的卒中风险因果效应是多少？\\
混杂&\chTenTablePath{2}&开放&关闭路径&在模型中调整吸烟&考虑吸烟的影响后，吸大麻对肺癌的因果效应是多少？\\
共同效应&\chTenTablePath{3}&关闭&保持关闭&不在模型中纳入体重下降&低热量饮食对胃癌的因果效应是多少？\\\bottomrule
\multicolumn{6}{p{154mm}}{\footnotesize $^{a}$涉及中介的分析有时会引入细微的复杂问题。如果中介与结局的关系有混杂因素，例如高血压可能混杂颈动脉狭窄与卒中的关联，纳入中介就可能打开一条原本关闭的路径，因为混杂因素和预测因素都指向中介。如果同时把该混杂因素纳入模型，就能再次关闭路径，因此不成问题。LDL为血清低密度脂蛋白胆固醇水平。}\\
\end{longtable}
\sourcepage{320}
对于涉及中介因素的路径，研究者可以根据目标作出选择。如果要确定预测因素对结局的总效应，就不应在分析模型中纳入中介因素。如果要确定预测因素的直接效应，即其不经中介因素发挥作用的部分，就应将中介因素，以及中介因素与结局关系的混杂因素，一并纳入模型。

\subsection{DAG与抽样}
DAG还可以澄清研究的抽样方式是否会引起选择偏倚。例如，接受绝经后激素替代治疗（hormone replacement therapy，HRT）的女性是否更容易患子宫内膜癌，曾经存在争论（36）。研究者面临的问题是如何处理绝经后子宫出血，因为它既是HRT的后果，也是子宫内膜癌的症状。

如果只将有绝经后出血的患者纳入样本，或只纳入没有绝经后出血的患者，就相当于把这一症状作为共同效应进行条件化。\sourcepage{321}HRT与子宫内膜癌都增加出血风险，因此只研究有出血者，或只研究无出血者，都可能得出错误结果（图10A.9）。只要预测因素和结局都影响某个参与者被纳入研究的可能性，就可能发生这一问题。
\par
\begin{figure}[htbp]\centering\chTenDAG{9}
\caption*{\textbf{图10A.9\quad 共同效应：绝经后出血碰撞点。}若影响入选，或被不恰当地纳入模型，就使HRT对子宫内膜癌因果估计偏倚。}
\end{figure}
\par

\subsection{共有原因}
本章前面提到，要成为混杂因素，一个变量必须引起所关注的暴露，或与该暴露具有共同原因，同时还引起结局。现在用DAG说明“具有共同原因”的含义。

考虑红肉摄入是否引起结肠癌的研究。结肠癌筛查是不是混杂因素？癌症筛查与结局存在因果联系：大多数结直肠癌病例通过筛查发现，有些甚至在不筛查时就不会被发现；但筛查大概不会使人吃或不吃红肉。不过，癌症筛查与不吃红肉具有一个共同原因，即采用“健康生活方式”。因此，从吃红肉经健康生活方式、癌症筛查到结直肠癌，存在一条需要关闭的混杂路径。即使无法直接测量健康生活方式本身，也可以通过调整癌症筛查将该路径关闭（图10A.10）。
\sourcepage{322}
\par
\begin{figure}[htbp]\centering\chTenDAG{10}
\caption*{\textbf{图10A.10\quad 共同原因。}健康生活方式使红肉经健康生活方式、癌症筛查至结直肠癌的路径开放；在模型纳入筛查可关闭。}
\end{figure}
\par
这一例子还说明了DAG的另一点。连接两个变量的箭头，方向始终应从因指向果，但它既可以表示增加，也可以表示降低该效应出现的可能性。本例中，健康生活方式被认为会降低吃红肉的可能性。有时，研究者在箭头旁用$+$号表示该原因增加效应出现的可能性，用$-$号表示降低这一可能性。

\subsection{综合例子}
上述内容相当抽象，连有经验的流行病学家也可能感到困惑。现在通过一个例子，回到红肉摄入是否增加结直肠癌风险，以及是否可能通过肠道微生物组发挥作用的问题。假设研究者可以取得数千名患者的电子健康记录。他们曾在健康维护项目中接受营养评估，项目还包括基线结肠镜，以及之后按推荐间隔进行的结肠镜检查。由于只纳入做过结肠镜者，前例中从红肉摄入经健康生活方式、癌症筛查到结直肠癌的路径就已经关闭了。自最初评估以来，现已有20多年的随访资料。此外，参与者的血液乃至粪便样本，也一直深冻保存。

研究者计划开展回顾性队列研究，纳入基线结肠镜正常的所有人。预测变量为营养师在基线时评估的患者每周红肉摄入量，结局变量为随后确诊结直肠癌。还将提取基线膳食纤维摄入量、期间接受粪便隐血试验（fecal occult blood testing，FOBT）的资料，以及这里暂不讨论的其他协变量。在一项巢式病例对照研究中，研究者将比较结直肠癌病例与未患癌者样本的粪便微生物组，同时考察与胃肠道癌症相关的多种遗传多态性。

\sourcepage{323}记住，研究者想知道吃红肉是否引起结直肠癌。因此，研究结束后的第一项分析，是考察吃红肉较多的参与者是否确实有更高的结直肠癌风险。不过，还需要决定如何考虑癌症相关基因多态性及膳食纤维摄入量等其他变量，因为这些变量可能造成红肉摄入与结肠癌之间的非因果关联。下面先分别讨论这些变量；图10A.11将它们画在同一张DAG中。
\par
\begin{figure}[htbp]\centering\chTenManyDAG
\caption*{\textbf{图10A.11\quad 包含多种关系的DAG。}估计吃红肉对结直肠癌的因果效应时，分析模型只需纳入有下划线的变量。政治观点和遗传多态性并不同时关联预测因素（红肉摄入）及结局（结直肠癌），无须纳入模型；不过，纳入与结直肠癌强相关的遗传多态性，可能提高模型精度。膳食纤维在这张DAG中是混杂因素，应纳入以关闭开放路径。微生物组的特定组成成分是中介因素，虚线下划线表示研究者必须决定是否纳入。FOBT阳性是共同效应，不应纳入分析模型，以免打开已经关闭的路径。}
\end{figure}
\par

\subsubsection{不同时与红肉和结直肠癌有因果联系的变量}
癌症相关基因的多态性可能影响癌症风险，却不太可能与一个人的饮食选择有关，因此估计红肉摄入对癌症的因果效应时，无须纳入这些变量。当然，研究者很容易通过分析来核实这些多态性是否与饮食无关。不过，如果某些多态性与结直肠癌有很强的关联，把它们纳入分析模型会提高精度，\sourcepage{324}也能考察这些多态性、红肉摄入及结直肠癌之间的交互作用。

同样，政治观点（假设确实测量了它）虽然可能与吃红肉有关，却不太可能影响患者患结直肠癌的风险；这一假设也很容易检验。

\subsubsection{同时关联两者的变量}
\textbf{微生物组组成。}研究问题提出了饮食可能影响癌症风险的一种解释，即饮食影响肠道微生物组的组成，例如提高脆弱拟杆菌（\textit{Bacteroides fragilis}）或具核梭杆菌（\textit{Fusobacterium nucleatum}）的比例。因此，这些组成成分是潜在中介因素。通过分析，可以确定微生物组信息是否影响红肉摄入与结直肠癌的关联，以及如何影响。是否在调整模型中纳入微生物资料，取决于研究者是想确定吃红肉对结直肠癌的总效应，还是只想确定通过或不通过微生物组传递的效应。

\textbf{膳食纤维摄入。}另一方面，它是潜在混杂因素。摄入大量膳食纤维的患者，如纯素食者和素食者，较少吃红肉；他们也可能由于这一原因或其他原因，结直肠癌风险较低。必须在模型中纳入纤维摄入量，关闭这条混杂路径。（注意，也可能是较高的红肉摄入导致膳食纤维摄入减少，因此箭头可以朝另一个方向。在这种情况下，将膳食纤维纳入分析模型，可确定红肉摄入对结直肠癌的直接效应。）

\textbf{粪便隐血试验（FOBT）。}它是共同效应：红肉摄入和癌症发生都可能导致试验阳性，并引出后续的进一步检查。这条路径已经关闭，因为两个箭头在FOBT处相遇。因此，不应在模型中纳入FOBT结果，以免把这条路径打开。

\subsection{DAG的局限}
使用DAG并不能完美解决因果推断中的所有问题。图形常常很容易画，但确定其背后的因果模型可能很困难：并非所有三变量关系都清楚，更不用说四、五或十个变量的关系。\sourcepage{325}有时，箭头方向并不明显，也不清楚两个变量是否应当相连。有时箭头似乎应是双向的：究竟是抑郁使人缺乏运动，还是反过来？我们还已经见过变量关系似乎形成循环的例子：肥胖引起关节炎，关节炎引起缺乏运动，缺乏运动又引起肥胖，如此循环。

考虑的变量越多，跟踪路径、判断变量间联系就越复杂。事实上，画DAG也能帮助识别哪些变量彼此无关，因此无须调整。最后，另一个局限是，DAG没有直接的方法来表示或处理交互作用及效应修饰。

好在设计研究时，无须完全确定最终DAG的样子；不过，为避免对共同效应条件化，必须确保潜在碰撞点不影响抽样计划。最终DAG可以等到分析结果、进行因果推断时再确定，但应当测量那些认为很可能出现在因果模型中的变量，即使尚不确定它们最终会发挥什么作用。

\chapterappendix{10B}{混杂与效应修饰的假想例子}
\sourcepage{326}
\definecolor{chTenRed}{HTML}{BF0040}
\definecolor{chTenPurple}{HTML}{633EA6}
\definecolor{chTenCyan}{HTML}{009FE3}
\definecolor{chTenTeal}{HTML}{008B81}
下表中的数字，是一项关于吃红肉是否为结直肠癌（colorectal cancer，CRC）危险因素的假想病例对照研究的参与者人数。CRC＋表示结直肠癌病例，CRC－表示对照；红肉＋表示该病例或对照吃红肉，红肉－表示不吃。部分数字以颜色标示，便于看出它们的来源。

\textbf{面板1。}左侧考察全体研究参与者时，吃红肉与CRC关联的比值比为2.25。右侧为指导分析的DAG。
\begin{center}
\begin{minipage}{.54\textwidth}\centering
\begin{tabular}{lrrr}\toprule 男女合计&CRC＋&CRC－&合计\\\midrule
红肉＋&90&60&\textcolor{chTenRed}{150}\\红肉－&60&90&\textcolor{chTenPurple}{150}\\合计&\textcolor{chTenCyan}{150}&\textcolor{chTenTeal}{150}&300\\\bottomrule
\end{tabular}\par\medskip $\mathrm{OR}=\dfrac{90\times90}{60\times60}=2.25$
\end{minipage}\hfill
\begin{minipage}{.4\textwidth}\centering\chTenSexDAG\end{minipage}
\end{center}

\textbf{面板2。}根据DAG，上述关联可能有一部分由混杂造成。按性别分层的表格如下：
\begin{center}\small
\begin{tabular}{lrrr}\toprule 男性&CRC＋&CRC－&合计\\\midrule
红肉＋&78&42&\textcolor{chTenRed}{120}\\红肉－&32&28&\textcolor{chTenPurple}{60}\\合计&\textcolor{chTenCyan}{110}&\textcolor{chTenTeal}{70}&180\\\bottomrule
\end{tabular}\quad
\begin{tabular}{lrrr}\toprule 女性&CRC＋&CRC－&合计\\\midrule
红肉＋&12&18&\textcolor{chTenRed}{30}\\红肉－&28&62&\textcolor{chTenPurple}{90}\\合计&\textcolor{chTenCyan}{40}&\textcolor{chTenTeal}{80}&120\\\bottomrule
\end{tabular}
\par\medskip $\mathrm{OR}_{\text{男}}=\dfrac{78\times28}{32\times42}=1.63$；\quad
$\mathrm{OR}_{\text{女}}=\dfrac{12\times62}{28\times18}=1.48$。
\end{center}
各层比值比明显低于2.25：男性为1.63，女性为1.48。可见性别既与CRC有关，病例中73\%（110／150）为男性，对照中仅47\%（70／150）为男性；也与吃红肉有关，男性中67\%（120／180）吃红肉，而女性仅25\%（30／120）吃红肉。\sourcepage{327}在本例中，按性别分层减弱了红肉与CRC的关联，但没有消除这一关联。

\textbf{面板3。效应修饰。}如果红肉摄入与CRC的关联在男性和女性中明显不同，面板1中的关联就受到性别的效应修饰。当然，像面板2那样，各层效应指标有小幅差别是可以预期的。但在以下表格中，交互作用很大：面板1中总比值比2.25，来自男性的比值比4.70，以及女性仅0.73的比值比。
\begin{center}\small
\begin{tabular}{lrrr}\toprule 男性&CRC＋&CRC－&合计\\\midrule
红肉＋&78&42&120\\红肉－&17&43&60\\合计&95&85&180\\\bottomrule
\end{tabular}\quad
\begin{tabular}{lrrr}\toprule 女性&CRC＋&CRC－&合计\\\midrule
红肉＋&12&18&30\\红肉－&43&47&90\\合计&55&65&120\\\bottomrule
\end{tabular}
\par\medskip $\mathrm{OR}_{\text{男}}=\dfrac{78\times43}{17\times42}=4.70$；\quad
$\mathrm{OR}_{\text{女}}=\dfrac{12\times47}{43\times18}=0.73$。
\end{center}
存在效应修饰时，各层比值比不同，应分别报告。

当然，实际研究资料中的数字，不像编造的数据那样泾渭分明。这时，可以结合研究以外的知识与统计学显著性检验，判断观察到的表面混杂或效应修饰程度，是否能合理地由资料中的偶然变异解释。

\chapterappendix{10C}{建模的简化例子}
\sourcepage{328}
假设一项研究发现，父母教育程度和儿童血铅水平是儿童IQ的两个主要预测因素。考虑以下铅水平正常与升高儿童的假想资料：
\begin{center}\begin{tabular}{lrr}\toprule&父母平均教育年限&儿童平均IQ\\\midrule
高铅&10.0&95\\正常铅&12.0&110\\\bottomrule\end{tabular}\end{center}
注意，父母教育程度也与儿童血铅水平有关。问题是：“铅水平正常与升高儿童之间的IQ差异，是否超过仅由父母教育程度差异所能解释的部分？”为回答这个问题，先考察父母教育程度的差异，预期会造成多大的IQ差异。方法是在铅水平正常的儿童中，绘制父母教育程度与儿童IQ的关系（图10C.1）。\footnote{这里简化了\term{协方差分析}（analysis of covariance，ANCOVA）的说明。实际上，应分别在正常铅组和高铅组绘制父母教育程度与儿童IQ的关系，再使用最适合两组图形的共同斜率。因此，这种调整模型假定两组中父母教育程度与IQ均为线性关系，且两组直线的斜率相同。其他更复杂的“回归调整”模型，则允许斜率和截距不同。}
\par
\begin{center}\begin{minipage}{\textwidth}\centering\chTenIQ
\captionof*{figure}{\textbf{图10C.1\quad 假想图：儿童IQ与父母教育年限之间的线性关系（虚线）。}}
\end{minipage}\end{center}
\par
\sourcepage{329}
图10C.1中的斜向虚线，表示铅水平正常儿童的IQ与父母教育程度之间的关系：父母每多接受两年教育，儿童IQ就增加5分。因此，可以从$A$点沿直线向下移动到$A'$点，调整正常铅组的IQ，以消除父母平均教育程度的差异。正常铅组的父母平均多接受了两年教育，所以把该组IQ下调5分，使其与高铅组在父母平均教育程度上具有可比性。这样，$A$与$B$两点之间仍有10分的IQ差，提示铅对IQ有这一大小的独立影响。\footnote{译注：底本此句写$A$与$B$相差10分；图中$A$为110、$B$为95，调整后10分差实际位于$A'$与$B$。}因此，低铅与高铅儿童原有的15分IQ差异中，5分可由父母教育程度不同解释，其余10分归于铅暴露。

\chapterappendix{10D}{第10章习题：用观察性研究估计因果效应}
\sourcepage{330}
\begin{enumerate}
\item 研究者开展一项病例对照研究，回答以下研究问题：“多吃水果和蔬菜能降低冠心病（coronary heart disease，CHD）风险吗？”假设研究发现，水果和蔬菜摄入量高于中位数者的比值比为0.6。
\begin{enumerate}
\item 观察到水果和蔬菜摄入较多者的CHD风险较低，可能有哪些原因使这一关联不代表因果效应？请特别考虑运动的混杂作用：如果多吃水果和蔬菜的人也运动更多，而运动正是他们CHD发生率较低的原因，就可能出现这种情况。
\item 可以采用哪些方法，应对运动这一潜在混杂因素？每种方案的优点和缺点是什么？
\end{enumerate}
\item 儿科诊所研究网络（Pediatric Research in Office Settings，PROS）的儿科医师开展的一项研究发现，在因发热被带到儿科医师处就诊的幼小婴儿（不足3个月）中，未接受包皮环切的男婴患尿路感染的风险，约为接受环切者的10倍（39）；这一关联已见于许多研究。有趣的是，在该研究中，未环切男婴的耳感染风险似乎较低（风险比0.77，$P=0.08$）。请解释：为什么只纳入发热婴儿，会在包皮环切与耳感染之间引入一种一般幼小婴儿人群中并不存在的关联？
\item 在第2章习题2中，我们请你提出研究方案，以回答对乙酰氨基酚是否引起哮喘。对此关联提出的一种机制是，对乙酰氨基酚消耗谷胱甘肽；后者能保护肺部免受可能促成哮喘的氧化损伤。请简要描述，如何利用母亲抗氧化相关基因型的变异，增强“母亲使用对乙酰氨基酚与子代哮喘的关联属于因果关系”这一推断。
\end{enumerate}
\sourcepage{331}
